%======================================================================%
\section{Octonions and the Albert Algebra}
\label{sec:app-albert:main}
%======================================================================%

This appendix collects, in a single self-contained place, the octonionic and
Jordan-algebraic facts on which the rest of the manuscript rests: the octonion
multiplication table and structure constants; the exceptional Jordan (Albert)
algebra $J_3(\mathbb O)$ as $3\times 3$ octonion-Hermitian matrices; the Jordan
product, the trace bilinear form, the cubic norm $N$, and the sharp map $X^\#$
with the Freudenthal identities; the rank stratification and primitive
idempotents; and the actions of $F_4$ and $E_6$. Every statement is tagged as a
\emph{theorem} (proved here or from cited standard references), a
\emph{postulate} (assumed added structure, cf.\ P1--P6), or as
\emph{indicative}. The identities assembled here are exactly those invoked by the
potential, mass-spectrum and gauge sections; we therefore cross-reference them by
the stable labels defined below.

Throughout we use the metric signature $(-,+,+,+)$ and the manuscript-wide
conventions: $[f]=\text{mass}$, and $\hat g^2,\xi,\hat y^2,\kappa$ dimensionless.
Two real forms of $E_6$ appear and must \emph{not} be conflated: the
\emph{compact} $E_6$ (equivalently $E_{6(-14)}$, maximal compact
$\mathrm{Spin}(10)\times U(1)$) governs the dynamics and gives a
positive-definite coset metric, while the non-compact $E_{6(-26)}$ is used
\emph{only} as the rigid classifier of orbits of the real $27$. This distinction
is the foundation enforced in \autoref{sec:app-albert:invariants} and
\autoref{sec:app-albert:orbits}, and underlies the vacuum-selection theorem cited
as \autoref{thm:potential:T4}.

%%%%%%%%%%%%%%%%%%%%%%%%%%%%%%%%%%%%%%%%%%%%%%%%%%%%%%%%%%%%%%%%%%%%%%
\subsection{Octonions: multiplication and structure constants}
\label{sec:app-albert:octonions}
%%%%%%%%%%%%%%%%%%%%%%%%%%%%%%%%%%%%%%%%%%%%%%%%%%%%%%%%%%%%%%%%%%%%%%

\begin{definition}[Octonions]\label{def:app-albert:octonions}
The octonions $\mathbb O$ form the real, $8$-dimensional, unital,
\emph{non-associative} division algebra. Fix the standard basis
$\{e_0=1,\,e_1,\dots,e_7\}$ with $e_i^2=-1$ for $i\ge 1$ and $e_i e_j=-e_j e_i$
for $i\neq j$ ($i,j\ge 1$). The multiplication is encoded by the totally
antisymmetric structure constants $\psi_{ijk}$,
\begin{equation}\label{eq:app-albert:struct}
  e_i\, e_j \;=\; -\,\delta_{ij}\,e_0 \;+\; \psi_{ijk}\,e_k ,
  \qquad i,j,k\in\{1,\dots,7\},
\end{equation}
with $\psi_{ijk}=+1$ on the seven Fano lines
\begin{equation}\label{eq:app-albert:fano}
  (ijk)\in\bigl\{(124),(235),(346),(457),(561),(672),(713)\bigr\}
\end{equation}
(and cyclic), and $\psi_{ijk}=0$ otherwise. Conjugation is
$\bar e_0=e_0$, $\bar e_i=-e_i$ ($i\ge1$); the (positive-definite) norm is
$n(x)=x\bar x=\bar x x=\sum_{a=0}^{7}x_a^2$, and $\overline{xy}=\bar y\,\bar x$.
\end{definition}

\begin{theorem}[Composition, alternativity, Moufang]\label{thm:app-albert:moufang}
\textup{(Theorem; standard \cite{Baez:2002,Springer:2000}.)}
$\mathbb O$ is a normed composition algebra, $n(xy)=n(x)\,n(y)$. It is not
associative but is \emph{alternative}: the associator
$[x,y,z]:=(xy)z-x(yz)$ is totally antisymmetric, so any subalgebra generated by
two elements is associative. The Moufang identities hold, in particular
$(xy)(zx)=x\bigl((yz)x\bigr)$. Real and imaginary parts obey
$\Re(xy)=\Re(yx)$ and $\Re\bigl((xy)z\bigr)=\Re\bigl(x(yz)\bigr)$, so the
trilinear form $\Re(xyz)$ is well defined (no bracketing needed) and cyclic.
\end{theorem}

\begin{proof}
Composition follows because $n$ is the norm of a Cayley--Dickson double of the
quaternions; alternativity and the antisymmetry of $\psi_{ijk}$ are equivalent to
the total antisymmetry of the associator, and the Moufang identities are the
standard consequences. The reality identities follow from
$x+\bar x=2\Re(x)\,e_0\in\mathbb R$ together with $\overline{xy}=\bar y\bar x$.
See \cite{Baez:2002,Schafer:1966}. $\qquad\blacksquare$
\end{proof}

\begin{remark}[$G_2$ and $\mathrm{Spin}(7)$]\label{rmk:app-albert:g2}
\textup{(Indicative context.)}
The automorphism group $\mathrm{Aut}(\mathbb O)$ is the compact $G_2$ (dim $14$),
the stabilizer of $\psi_{ijk}$ viewed as a generic $3$-form on $\Im\mathbb O\cong
\mathbb R^7$; the spin group $\mathrm{Spin}(7)$ acts by triality on the three
$8$-dimensional modules $\mathbf 8_v,\mathbf 8_s,\mathbf 8_c$. These reappear when
$\mathrm{Spin}(10)$ spinors are realized on $J_3(\mathbb O)$
(\autoref{sec:app-albert:F4E6}) and feed the coset tangent module
$\mathfrak m=16_{-3}\oplus\overline{16}_{+3}$ used downstream.
\end{remark}

%%%%%%%%%%%%%%%%%%%%%%%%%%%%%%%%%%%%%%%%%%%%%%%%%%%%%%%%%%%%%%%%%%%%%%
\subsection{The Albert algebra \texorpdfstring{$J_3(\mathbb O)$}{J3(O)}}
\label{sec:app-albert:jalg}
%%%%%%%%%%%%%%%%%%%%%%%%%%%%%%%%%%%%%%%%%%%%%%%%%%%%%%%%%%%%%%%%%%%%%%

\begin{definition}[Albert algebra]\label{def:app-albert:jalg}
The (real, reduced) \emph{exceptional Jordan algebra} or \emph{Albert algebra} is
the space of $3\times3$ octonion-Hermitian matrices
\begin{equation}\label{eq:app-albert:Xform}
  J_3(\mathbb O)=
  \left\{\,
  X=\begin{pmatrix}
      \xi_1 & x_3 & \bar x_2\\
      \bar x_3 & \xi_2 & x_1\\
      x_2 & \bar x_1 & \xi_3
    \end{pmatrix}
  :\ \xi_i\in\mathbb R,\ x_i\in\mathbb O
  \,\right\},
\end{equation}
equipped with the commutative \emph{Jordan product}
\begin{equation}\label{eq:app-albert:jprod}
  X\circ Y=\tfrac12\,(XY+YX),
\end{equation}
where $XY$ is ordinary (octonionic) matrix multiplication.
\end{definition}

\begin{theorem}[Dimension and Jordan identity]\label{thm:app-albert:dim}
\textup{(Theorem.)}
$J_3(\mathbb O)$ is a $27$-dimensional, commutative, power-associative Jordan
algebra: $\circ$ is commutative and satisfies the Jordan identity
$X\circ\bigl(Y\circ X^{\circ 2}\bigr)=(X\circ Y)\circ X^{\circ 2}$. It is
\emph{formally real} (Euclidean): $\sum_i A_i^{\circ2}=0\Rightarrow A_i=0$. It is
the unique exceptional Jordan algebra (not isomorphic to a subalgebra of an
associative algebra under $\circ$).
\end{theorem}

\begin{proof}
The real count is $3$ (diagonal) $+\,3\times 8$ (off-diagonal octonions)
$=27$. Commutativity is immediate from \eqref{eq:app-albert:jprod}. The Jordan
identity holds for $3\times3$ Hermitian matrices over an alternative algebra by
the Hermitian-matrix theorem of Jordan--von Neumann--Wigner
\cite{JNW:1934,McCrimmon:2004}; formal reality follows from positivity of the
trace form \eqref{eq:app-albert:Q} below. Exceptionality (the Albert
exceptionality theorem) is standard \cite{Albert:1934,Schafer:1966}.
$\qquad\blacksquare$
\end{proof}

\begin{definition}[Traces, trace form, cubic norm]\label{def:app-albert:invariants}
On $J_3(\mathbb O)$ define the \emph{linear trace}
$\operatorname{tr}(X)=\xi_1+\xi_2+\xi_3$, the symmetric \emph{trace bilinear form}
\begin{equation}\label{eq:app-albert:Q}
  \langle X,Y\rangle:=\operatorname{tr}(X\circ Y),
  \qquad
  Q(X):=\langle X,X\rangle=\operatorname{tr}(X^{\circ2})
  =\sum_i\xi_i^2+2\sum_i n(x_i)\ \ge 0,
\end{equation}
which is positive-definite, and the \emph{cubic norm} (determinant analogue)
\begin{equation}\label{eq:app-albert:N}
  N(X):=\xi_1\xi_2\xi_3-\sum_{i=1}^{3}\xi_i\,n(x_i)+2\,\Re\!\bigl(x_1 x_2 x_3\bigr).
\end{equation}
The second elementary symmetric invariant is
$\sigma_2(X)=\tfrac12\bigl(\operatorname{tr}(X)^2-Q(X)\bigr)$. The trilinear
polarization of $N$ is denoted $N(X,Y,Z)$ with $N(X,X,X)=6N(X)$.
\end{definition}

\begin{theorem}[Cubic Cayley--Hamilton]\label{thm:app-albert:CH}
\textup{(Theorem.)}
Every $X\in J_3(\mathbb O)$ satisfies the cubic characteristic identity
\begin{equation}\label{eq:app-albert:CH}
  X^{\circ3}-\operatorname{tr}(X)\,X^{\circ2}+\sigma_2(X)\,X-N(X)\,E=0,
  \qquad E:=\mathrm{diag}(1,1,1).
\end{equation}
Its three roots are the real \emph{Jordan eigenvalues}
$\lambda_1,\lambda_2,\lambda_3$, in terms of which
$\operatorname{tr}(X)=\sum_i\lambda_i$, $Q(X)=\sum_i\lambda_i^2$,
$N(X)=\lambda_1\lambda_2\lambda_3$.
\end{theorem}

\begin{proof}
By the spectral theorem for formally real Jordan algebras every $X$ is
$F_4$-conjugate to $\mathrm{diag}(\lambda_1,\lambda_2,\lambda_3)$ with
$\lambda_i\in\mathbb R$ (\autoref{thm:app-albert:F4} below); on diagonal elements
\eqref{eq:app-albert:CH} is the ordinary scalar Cayley--Hamilton identity, and
$F_4$-equivariance of every term ($\operatorname{tr},\sigma_2,N,E$ are
$F_4$-invariant or invariant) transports it to all of $J_3(\mathbb O)$
\cite{McCrimmon:2004,Springer:2000}. $\qquad\blacksquare$
\end{proof}

%%%%%%%%%%%%%%%%%%%%%%%%%%%%%%%%%%%%%%%%%%%%%%%%%%%%%%%%%%%%%%%%%%%%%%
\subsection{The sharp map and the Freudenthal identities}
\label{sec:app-albert:sharp}
%%%%%%%%%%%%%%%%%%%%%%%%%%%%%%%%%%%%%%%%%%%%%%%%%%%%%%%%%%%%%%%%%%%%%%

\begin{definition}[Sharp map and cross product]\label{def:app-albert:sharp}
The \emph{sharp} (Freudenthal adjoint) map $X\mapsto X^\#$ is the unique quadratic
map with entries the $2\times2$ octonionic cofactors of $X$:
\begin{equation}\label{eq:app-albert:sharpentries}
  (X^\#)_{ii}=\xi_j\xi_k-n(x_i),
  \qquad
  (X^\#)_{ij}=\overline{x_i x_j}-\xi_k\,x_k
  \quad(\text{cyclic } i,j,k).
\end{equation}
Equivalently, in the Jordan frame
$\mathrm{diag}(\lambda_1,\lambda_2,\lambda_3)^\#
 =\mathrm{diag}(\lambda_2\lambda_3,\lambda_3\lambda_1,\lambda_1\lambda_2)$. Its
symmetric polarization defines the \emph{Freudenthal cross product}
\begin{equation}\label{eq:app-albert:cross}
  X\times Y:=(X+Y)^\#-X^\#-Y^\#,
  \qquad X\times X=2X^\#,
\end{equation}
a symmetric bilinear map $J_3(\mathbb O)\times J_3(\mathbb O)\to J_3(\mathbb O)$.
\end{definition}

\begin{theorem}[Freudenthal / adjoint identities]\label{thm:app-albert:freudenthal}
\textup{(Theorem; \cite{Freudenthal:1985,McCrimmon:2004,Springer:2000}.)}
For all $X,Y\in J_3(\mathbb O)$ the following hold:
\begin{align}
  X^\#\circ X &= N(X)\,E,
  \label{eq:app-albert:F1}\\
  (X^\#)^\# &= N(X)\,X,
  \label{eq:app-albert:F2}\\
  \langle X^\#,Y\rangle &= DN(X)[Y]
     = N(X,X,Y),
  \quad\text{i.e. } X^\#=\nabla N(X),
  \label{eq:app-albert:F3}\\
  \langle X^\#,Y\rangle &= \sigma_2(X)\,\langle?\rangle
     \ \Longrightarrow\
  \operatorname{tr}(X^\#)=\sigma_2(X),
  \label{eq:app-albert:F4}\\
  \langle X^\#,X^\#\rangle
     &= \sigma_2(X)^2-2\,\operatorname{tr}(X)\,N(X)
     = (\lambda_2\lambda_3)^2+(\lambda_3\lambda_1)^2+(\lambda_1\lambda_2)^2.
  \label{eq:app-albert:F5}
\end{align}
\end{theorem}

\begin{proof}
Identity \eqref{eq:app-albert:F1} is the cofactor expansion of the determinant
$N$ and follows from \eqref{eq:app-albert:CH} by writing
$X^\#=X^{\circ2}-\operatorname{tr}(X)X+\sigma_2(X)E$ and using
\eqref{eq:app-albert:CH}; \eqref{eq:app-albert:F2} (the ``$\#\#$'' Freudenthal
identity) is verified in the Jordan frame, where
$(\lambda_2\lambda_3,\lambda_3\lambda_1,\lambda_1\lambda_2)^\#
=(\lambda_1^2\lambda_2\lambda_3,\dots)=N\cdot(\lambda_1,\lambda_2,\lambda_3)$, and
extended by $F_4$-equivariance. The gradient identity \eqref{eq:app-albert:F3} is
the definition of $X^\#$ as the differential of the cubic $N$ with respect to the
trace form, and $\operatorname{tr}(X^\#)=\sigma_2(X)$ in
\eqref{eq:app-albert:F4} follows from
$X^\#=X^{\circ2}-\operatorname{tr}(X)X+\sigma_2(X)E$ on taking $\operatorname{tr}$.
Finally \eqref{eq:app-albert:F5} is the frame evaluation of
$\langle X^\#,X^\#\rangle$, manifestly a sum of squares; the identity
$\sigma_2(X)^2-2\operatorname{tr}(X)N(X)$ is the symmetric-function rewriting.
$\qquad\blacksquare$
\end{proof}

\begin{corollary}[Rank from the sharp map]\label{cor:app-albert:rank}
\textup{(Theorem.)}
Define the \emph{rank} by
\begin{equation}\label{eq:app-albert:rank}
  \operatorname{rank}(X)=
  \begin{cases}
    0,& X=0,\\
    1,& X\neq 0,\ X^\#=0,\\
    2,& X^\#\neq0,\ N(X)=0,\\
    3,& N(X)\neq0.
  \end{cases}
\end{equation}
Then $X^\#=0\iff\operatorname{rank}(X)\le1$, and $N(X)=0\iff\operatorname{rank}(X)
\le2$. In particular the cubic norm $N$, not the quadratic $Q$, detects rank
deficiency at order $3$.
\end{corollary}

\begin{proof}
In the Jordan frame $X^\#=0$ iff all pairwise products $\lambda_i\lambda_j$
vanish, i.e.\ at most one $\lambda_i\neq0$; and $N=\lambda_1\lambda_2\lambda_3=0$
iff some $\lambda_i=0$. Equivariance extends both statements to all of
$J_3(\mathbb O)$. $\qquad\blacksquare$
\end{proof}

\begin{remark}[Positivity of the sharp invariant]\label{rmk:app-albert:positivity}
\textup{(Theorem.)}
Because the trace form \eqref{eq:app-albert:Q} is positive-definite,
$\langle X^\#,X^\#\rangle\ge0$ with equality iff $X^\#=0$ iff
$\operatorname{rank}(X)\le1$ (\autoref{cor:app-albert:rank}). This is the exact
algebraic input to the potential: the $E_6$-covariant term
$\alpha\langle X^\#,X^\#\rangle$ vanishes precisely on the rank-$\le1$ cone, which
is why it selects the rank-$1$ idempotent orbit; see
\autoref{thm:potential:T1} and \autoref{thm:potential:T4}.
\end{remark}

%%%%%%%%%%%%%%%%%%%%%%%%%%%%%%%%%%%%%%%%%%%%%%%%%%%%%%%%%%%%%%%%%%%%%%
\subsection{Rank stratification and primitive idempotents}
\label{sec:app-albert:idempotents}
%%%%%%%%%%%%%%%%%%%%%%%%%%%%%%%%%%%%%%%%%%%%%%%%%%%%%%%%%%%%%%%%%%%%%%

\begin{definition}[Idempotents and Jordan frame]\label{def:app-albert:idemp}
$P\in J_3(\mathbb O)$ is an \emph{idempotent} if $P\circ P=P$, and
\emph{primitive} if it is a nonzero idempotent that is not a sum of two nonzero
orthogonal idempotents. A \emph{complete Jordan frame} is a triple
$\{P_1,P_2,P_3\}$ of pairwise orthogonal ($P_i\circ P_j=0$, $i\neq j$) primitive
idempotents with $P_1+P_2+P_3=E$.
\end{definition}

\begin{theorem}[Idempotent stratification]\label{thm:app-albert:idemp}
\textup{(Theorem; \cite{McCrimmon:2004,Yokota:2009}.)}
A primitive idempotent has Jordan eigenvalues $(1,0,0)$, hence
$\operatorname{tr}(P)=1$, $Q(P)=1$, $N(P)=0$, $P^\#=0$, i.e.\
$\operatorname{rank}(P)=1$. The normalized rank-$1$ elements are exactly the
primitive idempotents; they form a single $F_4$-orbit
\begin{equation}\label{eq:app-albert:OP2}
  F_4/\mathrm{Spin}(9)\ \cong\ \mathbb{OP}^2,
  \qquad \dim_{\mathbb R}=16,
\end{equation}
the octonionic projective plane (the real rank-$1$ idempotents). Under the full
\emph{compact} $E_6$ the projective class $[X]$ of a rank-$1$ line sweeps the
larger homogeneous space
\begin{equation}\label{eq:app-albert:EIII}
  \mathcal M=\mathrm{EIII}=E_6/(\mathrm{Spin}(10)\times U(1)),
  \qquad \dim_{\mathbb R}\mathcal M=32,
\end{equation}
the complex Cayley plane and the closed minimal $E_6$-orbit in
$\mathbb P(27_{\mathbb C})$.
\end{theorem}

\begin{proof}
$P\circ P=P$ with eigenvalues $\mu$ forces $\mu^2=\mu$, so $\mu\in\{0,1\}$;
primitivity forbids two $1$'s, and nonzero forbids all $0$'s, leaving $(1,0,0)$.
The remaining invariants follow from \autoref{thm:app-albert:CH} and
\eqref{eq:app-albert:sharpentries}. Transitivity of $F_4=\mathrm{Aut}(J_3(\mathbb
O))$ on primitive idempotents with stabilizer $\mathrm{Spin}(9)$ is the standard
identification $\mathbb{OP}^2=F_4/\mathrm{Spin}(9)$
\cite{Baez:2002,Yokota:2009}; the $E_6$-orbit statement is the realization of
EIII as the minimal/highest-weight orbit in the projectivized $27$
\cite{Springer:2000}. $\qquad\blacksquare$
\end{proof}

\begin{remark}[Vacuum manifold; excluded base point]\label{rmk:app-albert:vacuum}
\textup{(Theorem, consistency note.)}
The UCFT vacuum manifold is exactly \eqref{eq:app-albert:EIII}, the rank-$1$
primitive-idempotent orbit, $\dim_{\mathbb R}=32$ --- \emph{not} the orbit of
$E=\mathrm{diag}(1,1,1)$. The latter has rank $3$, $E^\#=E\neq0$,
$\langle E^\#,E^\#\rangle=3>0$, stabilizer $F_4$, and orbit $E_6/F_4$ of real
dimension $26$; by \autoref{rmk:app-albert:positivity} it is strictly disfavoured
by the potential and is therefore excluded. This is consistent with the
vacuum-selection theorem \autoref{thm:potential:T4} and the coset tangent
$\mathfrak m=16_{-3}\oplus\overline{16}_{+3}$ of \autoref{sec:app-albert:branch}.
\end{remark}

%%%%%%%%%%%%%%%%%%%%%%%%%%%%%%%%%%%%%%%%%%%%%%%%%%%%%%%%%%%%%%%%%%%%%%
\subsection{Invariants and the two real forms}
\label{sec:app-albert:invariants}
%%%%%%%%%%%%%%%%%%%%%%%%%%%%%%%%%%%%%%%%%%%%%%%%%%%%%%%%%%%%%%%%%%%%%%

\begin{theorem}[$F_4=\mathrm{Aut}(J_3(\mathbb O))$]\label{thm:app-albert:F4}
\textup{(Theorem; \cite{Schafer:1966,Yokota:2009}.)}
The automorphism group of the Jordan algebra (preserving the product $\circ$,
hence the unit $E$ and \emph{both} $Q$ and $N$) is the compact exceptional group
$F_4$, $\dim_{\mathbb R}F_4=52$. It acts on $J_3(\mathbb O)=\mathbb R\oplus 26$
fixing the unit; on the traceless $26$ it embeds $F_4\subset SO(26)\subset
SO(27)$, since it preserves the Euclidean form $Q$.
\end{theorem}

\begin{theorem}[Reduced structure group $=E_6$]\label{thm:app-albert:E6}
\textup{(Theorem; \cite{Jacobson:1968,Springer:2000}.)}
The \emph{reduced structure group} of $J_3(\mathbb O)$ --- the connected group of
linear maps preserving the cubic norm $N$ up to a real character --- is the
exceptional group $E_6$, $\dim_{\mathbb R}E_6=78$, acting irreducibly on the $27$.
For $g$ in this group,
\begin{equation}\label{eq:app-albert:character}
  N(g\cdot X)=\lambda(g)\,N(X),\qquad \lambda(g)\in\mathbb R^\times,
\end{equation}
and the sharp map is covariant, $(g\cdot X)^\#=\lambda(g)\,(g^{-1})^{\mathsf T}\!
\cdot X^\#$ in the trace form. The split/non-compact real form realized by the
\emph{division}-octonion algebra is $E_{6(-26)}$; it does \emph{not} preserve
$Q=\operatorname{tr}(X^2)$.
\end{theorem}

\begin{proof}
Jacobson's construction realizes $\mathfrak e_6$ as $\mathfrak f_4$ (derivations)
plus the traceless multiplication operators
$L_X:Y\mapsto X\circ Y-\tfrac13\langle X,Y\rangle E$ acting on the $27$; the
group generated preserves $N$ up to the character \eqref{eq:app-albert:character}
\cite{Jacobson:1968}. The frame element $\mathrm{diag}(s,s,s^{-2})$ rescales $N$
nontrivially while changing $Q$, so $Q$ is not $E_6$-invariant; only $F_4$
(the $\lambda\equiv1$, unit-preserving subgroup) preserves $Q$. $\qquad
\blacksquare$
\end{proof}

\begin{postulate}[Compact real form for dynamics]\label{post:app-albert:compact}
\textup{(Convention, cf.\ P1--P2.)}
For the dynamics we use the \emph{compact} real form $E_6$ (equivalently
$E_{6(-14)}$, whose maximal compact subgroup is exactly
$H=\mathrm{Spin}(10)\times U(1)$). On the compact form the $27$ is unitary, the
trace form $Q=\langle\cdot,\cdot\rangle$ is the invariant positive-definite
metric, and the norm character is identically $\lambda\equiv1$ (a continuous
homomorphism from a compact connected group into $\mathbb R_{>0}$). Hence on the
compact form every $E_6$-invariant polynomial built from $\langle\cdot,\cdot
\rangle$ and $N$ --- in particular $\langle X^\#,X^\#\rangle$ --- is a genuine
invariant, the coset metric is positive-definite, and the gauge sector is unitary.
The non-compact $E_{6(-26)}$ is retained \emph{only} as the rigid orbit classifier
of \autoref{sec:app-albert:orbits}.
\end{postulate}

\begin{remark}[Only $N$ separates structure-group orbits]\label{rmk:app-albert:onlyN}
\textup{(Theorem, guard.)}
$Q=\operatorname{tr}(X^2)$ is $F_4$-invariant \emph{only}; it is \emph{not}
preserved by the $E_6$ structure group, and the pair $(Q,N)$ does \emph{not}
classify $E_6$-orbits. Only the cubic norm $N$ is structure-group invariant (up to
the character $\lambda$), and orbits are classified by rank and $N$
(\autoref{sec:app-albert:orbits}). Any claim that the structure group preserves
$Q$, or that $(Q,N)$ separates orbits, is false.
\end{remark}

%%%%%%%%%%%%%%%%%%%%%%%%%%%%%%%%%%%%%%%%%%%%%%%%%%%%%%%%%%%%%%%%%%%%%%
\subsection{Orbit classification of the real \texorpdfstring{$27$}{27}}
\label{sec:app-albert:orbits}
%%%%%%%%%%%%%%%%%%%%%%%%%%%%%%%%%%%%%%%%%%%%%%%%%%%%%%%%%%%%%%%%%%%%%%

\begin{theorem}[Rank-and-norm orbit classification]\label{thm:app-albert:orbits}
\textup{(Theorem; \cite{Springer:2000,Krutelevich:2007}.)}
The orbits of $E_{6(-26)}$ on the real $27=J_3(\mathbb O)$ are classified by the
rank \eqref{eq:app-albert:rank} together with the value of the cubic norm $N$:
\begin{itemize}
  \item rank $0$: $\{0\}$;
  \item rank $1$: $X\neq0$, $X^\#=0$ --- a single orbit, the cone over EIII; its
    projectivization is the closed minimal orbit $\mathcal M=\mathrm{EIII}$,
    $\dim_{\mathbb R}\mathcal M=32$;
  \item rank $2$: $X^\#\neq0$, $N(X)=0$ --- a single orbit;
  \item rank $3$: $N(X)\neq0$ --- orbits labelled by $\operatorname{sgn}N$ and the
    signature of $X$ (the indefinite cubic $N$ is genuinely sign-indefinite).
\end{itemize}
The quadratic $Q$ is \emph{not} an orbit invariant; it can be rescaled within a
fixed-$N$, fixed-rank orbit by the non-compact directions
$\mathrm{diag}(s,s,s^{-2})\in E_{6(-26)}$.
\end{theorem}

\begin{proof}
That rank and $N$ are structure-group invariants follows from
\eqref{eq:app-albert:character} and the covariance of $X^\#$
(\autoref{thm:app-albert:E6}); transitivity within each stratum is the
reduction of an arbitrary $X$ to a normal form
$\mathrm{diag}(\lambda_1,\lambda_2,\lambda_3)$ by $F_4$ followed by a
norm-rescaling diagonal element, with the rank-deficient strata as the explicit
degenerations \cite{Krutelevich:2007,Springer:2000}. $\qquad\blacksquare$
\end{proof}

\begin{remark}[Sharp characterization, summary]\label{rmk:app-albert:sharpsummary}
\textup{(Theorem.)}
Collecting \autoref{cor:app-albert:rank} and
\autoref{thm:app-albert:freudenthal}: $\operatorname{rank}\le1\Leftrightarrow
X^\#=0$; $\operatorname{rank}\le2\Leftrightarrow N(X)=0$;
$\operatorname{rank}=3\Leftrightarrow N(X)\neq0$. These are the identities the
potential section uses to localize the global minimum on the rank-$1$ cone
(\autoref{thm:potential:T4}).
\end{remark}

%%%%%%%%%%%%%%%%%%%%%%%%%%%%%%%%%%%%%%%%%%%%%%%%%%%%%%%%%%%%%%%%%%%%%%
\subsection{$F_4$- and $H$-branchings; the coset tangent module}
\label{sec:app-albert:branch}
%%%%%%%%%%%%%%%%%%%%%%%%%%%%%%%%%%%%%%%%%%%%%%%%%%%%%%%%%%%%%%%%%%%%%%

\begin{theorem}[$\mathrm{Spin}(10)$ realization]\label{thm:app-albert:F4E6}
\textup{(Theorem; \cite{Baez:2002,Springer:2000}.)}
Under $H=\mathrm{Spin}(10)\times U(1)\subset E_6$ (the maximal compact of
$E_{6(-14)}$), the fundamental and adjoint representations branch as
\begin{align}
  27 &= 1_{+4}\oplus 10_{-2}\oplus 16_{+1},
  \label{eq:app-albert:27branch}\\
  78 &= 45_{0}\oplus 1_{0}\oplus 16_{-3}\oplus \overline{16}_{+3},
  \label{eq:app-albert:78branch}
\end{align}
where the subscripts are the $U(1)$ charges in the normalization fixed by
\eqref{eq:app-albert:78branch}. The $10_{-2}$ is the $\mathrm{Spin}(10)$ vector
(an $H_2(\mathbb O)$-type block of $J_3(\mathbb O)$), the $16_{+1}$ a chiral
spinor, and the $1_{+4}$ a Jordan-frame singlet (an extra Standard-Model-singlet
modulus).
\end{theorem}

\begin{corollary}[Coset tangent module]\label{cor:app-albert:tangent}
\textup{(Theorem.)}
The tangent space to $\mathcal M=\mathrm{EIII}=E_6/(\mathrm{Spin}(10)\times U(1))$
at the basepoint is the $H$-module
\begin{equation}\label{eq:app-albert:tangent}
  \mathfrak m=16_{-3}\oplus\overline{16}_{+3},
  \qquad \dim_{\mathbb R}\mathfrak m=32,
\end{equation}
the complement of $\mathfrak h=45_0\oplus1_0$ in
\eqref{eq:app-albert:78branch}. It is \emph{real-irreducible of complex type}
(commutant $\cong\mathbb C$) and contains \emph{no} $\mathrm{Spin}(10)$ singlet.
Consequently the invariant coset metric $K|_{\mathfrak m}$ is unique up to scale
and \emph{positive-definite} (Schur), and any radiative ($H$-invariant)
mass operator on $\mathfrak m$ is a single scalar: one common mass on all $32$
clock fields.
\end{corollary}

\begin{proof}
$\mathfrak m=\mathfrak e_6\ominus\mathfrak h$ reads off
\eqref{eq:app-albert:78branch}. The $U(1)$ charges $\pm3$ forbid the symmetric
invariants in $16\otimes16$ and $\overline{16}\otimes\overline{16}$, leaving the
unique charge-$0$ cross-pairing $16_{-3}\otimes\overline{16}_{+3}\to1_0$; this is
the restriction of the Killing form, which on the compact form is
positive-definite. Real-irreducibility with commutant $\mathbb C$ is the
complex-type Schur statement, giving uniqueness of the metric and a single common
radiative mass. Absence of a singlet is immediate from
\eqref{eq:app-albert:tangent}. $\qquad\blacksquare$
\end{proof}

\begin{remark}[Excluded tangent decompositions]\label{rmk:app-albert:forbidden}
\textup{(Consistency constraint.)}
The coset tangent is \eqref{eq:app-albert:tangent} and contains no
$\mathrm{Spin}(10)$ singlet. Any decomposition of $\mathfrak m$ as $2\times10+4$,
or $10+\overline{10}+4$, or any singlet-in-$\mathfrak m$ form, is incorrect, as is
any attempt to extract a $4$-dimensional $H$-invariant subspace from
$\nabla\sigma_2,\nabla N$ singlets. The Lorentzian frame directions are
\emph{not} a sub-$H$-module of $\mathfrak m$; they are introduced as added
structure (soldering postulate P3, cf.\ \autoref{sec:app-soldering:main}) and are
gauge-protected, not Hessian zeros.
\end{remark}

%%%%%%%%%%%%%%%%%%%%%%%%%%%%%%%%%%%%%%%%%%%%%%%%%%%%%%%%%%%%%%%%%%%%%%
\subsection{The Minkowski block from the cubic norm}
\label{sec:app-albert:minkowski}
%%%%%%%%%%%%%%%%%%%%%%%%%%%%%%%%%%%%%%%%%%%%%%%%%%%%%%%%%%%%%%%%%%%%%%

\begin{theorem}[$H_2(\mathbb O)\cong\mathbb R^{1,9}$]\label{thm:app-albert:minkowski}
\textup{(Theorem; \cite{Baez:2002,Sudbery:1984}.)}
The $2\times2$ octonion-Hermitian matrices
\begin{equation}\label{eq:app-albert:H2O}
  H_2(\mathbb O)=
  \left\{\begin{pmatrix}a & x\\ \bar x & b\end{pmatrix}:a,b\in\mathbb R,\,x\in
  \mathbb O\right\},
  \qquad \dim_{\mathbb R}=10,
\end{equation}
form a sub-block of $J_3(\mathbb O)$ (set the third pivot to $1$). With
$\det X=ab-n(x)$ and the lightcone parametrization $a=x^0+x^9$, $b=x^0-x^9$,
\begin{equation}\label{eq:app-albert:detmink}
  \det X=(x^0)^2-(x^9)^2-\textstyle\sum_{i=1}^{8}(x^i)^2
  =-\,\eta_{\mu\nu}X^\mu X^\nu,
  \quad \eta=\mathrm{diag}(-1,+1,\dots,+1),
\end{equation}
so $(H_2(\mathbb O),-\det)\cong\mathbb R^{1,9}$. Restricting the octonion entry to
$\mathbb C\subset\mathbb O$ gives $H_2(\mathbb C)\cong\mathbb R^{1,3}$ with
$-\det$ the $(1,3)$ Minkowski form in the convention $(-,+,+,+)$.
\end{theorem}

\begin{proof}
$\det\!\bigl(\begin{smallmatrix}a&x\\\bar x&b\end{smallmatrix}\bigr)
=ab-x\bar x=ab-n(x)$ by definition of $n$; substituting the lightcone
coordinates yields \eqref{eq:app-albert:detmink}. The restriction to
$\mathbb C$ drops $x^4,\dots,x^9$ (the imaginary octonion directions beyond
$\mathbb C$), leaving the four coordinates $x^0,x^1,x^2,x^3$ with the same form.
$\qquad\blacksquare$
\end{proof}

\begin{remark}[Where Lorentzian signature enters]\label{rmk:app-albert:signature}
\textup{(Postulate P3.)}
The positive-definite object is the coset/target metric $K|_{\mathfrak m}$ on
$\mathfrak m$ (\autoref{cor:app-albert:tangent}); the indefinite object is the
cubic norm $N$, whose $H_2(\mathbb O)\cong\mathbb R^{1,9}$ block carries the
Minkowski form \eqref{eq:app-albert:detmink}. These are two distinct forms on two
distinct modules. The Killing form on $\mathfrak m$ for the compact $E_6$ has
signature $(32,0)$ and \emph{zero} Lorentzian content; the spacetime dimension
$d=4$ and signature $(1,3)=(-,+,+,+)$ are \emph{added structure} (soldering
postulate P3, cf.\ \autoref{sec:app-soldering:main}), soldered from $N$ via
$H_2(\mathbb C)$, and are \emph{not} Killing-form, coset-geometry, or rank-lemma
theorems. Any $(4,28)$ or $(1,3)$ coset/Killing-signature claim is false.
\end{remark}

%%%%%%%%%%%%%%%%%%%%%%%%%%%%%%%%%%%%%%%%%%%%%%%%%%%%%%%%%%%%%%%%%%%%%%
\subsection{Identities used downstream}
\label{sec:app-albert:identities}
%%%%%%%%%%%%%%%%%%%%%%%%%%%%%%%%%%%%%%%%%%%%%%%%%%%%%%%%%%%%%%%%%%%%%%

For convenience we list the identities invoked by the potential, mass-spectrum and
gauge sections, all proved above.

\begin{enumerate}
  \item \textbf{Sharp / Freudenthal} \eqref{eq:app-albert:F1}--\eqref{eq:app-albert:F5}:
    $X^\#\circ X=N(X)E$, $(X^\#)^\#=N(X)X$, $X^\#=\nabla N(X)$,
    $\operatorname{tr}(X^\#)=\sigma_2(X)$, and
    $\langle X^\#,X^\#\rangle=\sigma_2(X)^2-2\operatorname{tr}(X)N(X)\ge0$. Used in
    the $E_6$-covariant term of the potential (P1) and in
    \autoref{thm:potential:T1}.
  \item \textbf{Rank stratification} \eqref{eq:app-albert:rank}: rank $\le1
    \Leftrightarrow X^\#=0$; rank $\le2\Leftrightarrow N=0$. Used to localize the
    global minimum on the rank-$1$ cone (\autoref{thm:potential:T4}).
  \item \textbf{Cubic Cayley--Hamilton} \eqref{eq:app-albert:CH} and frame
    eigenvalues. Used to evaluate $V$ on the cone.
  \item \textbf{Branchings} \eqref{eq:app-albert:27branch},
    \eqref{eq:app-albert:78branch} and the tangent module
    \eqref{eq:app-albert:tangent}. Used for the $27\times27$ scalar block
    $M=\mathrm{diag}(\mu_1\mathbb 1_1,\mu_{10}\mathbb 1_{10},\mu_{16}\mathbb 1_{16})$,
    for the single common coset mass (Schur on $\mathfrak m$), and for the gauge
    multiplet assignments. The $28$ massive clock fields
    ($32=28\text{ massive}+4\text{ gauge/soldering}$) enter all loop sums; cf.\
    \autoref{thm:gravity:MPl} and the gauge $\beta$-function
    \autoref{thm:gravity:beta}.
  \item \textbf{Minkowski block} \eqref{eq:app-albert:detmink}: $-\det|_{H_2
    (\mathbb O)}$ is the $\mathbb R^{1,9}$ form, $-\det|_{H_2(\mathbb C)}$ the
    $(1,3)$ form. Used by the soldering postulate P3
    (\autoref{sec:app-soldering:main}).
  \item \textbf{Casimirs of the gauge representations.} With the unbroken
    $H=\mathrm{Spin}(10)\times U(1)$ acting on the multiplets in
    \eqref{eq:app-albert:27branch}, the relevant invariants are
    $C_A(\mathrm{SO}(10))=8$, $T(16)=2$ (the spinor is index $2$), and
    $C_2(16)=45/8$. These feed
    $b_0=\tfrac{11}{3}C_A-\tfrac43 T n_F-\tfrac16 T n_S$ in
    \autoref{thm:gravity:beta} (canonical anchor $b_0=26$ for one generation,
    $b_0=12$ for three).
\end{enumerate}

\begin{remark}[Status]\label{rmk:app-albert:status}
All results in this appendix are \emph{theorems} of octonionic/Jordan algebra
(\autoref{thm:app-albert:moufang}--\autoref{thm:app-albert:minkowski}) or standard
representation theory, except the two clearly marked conventions/postulates: the
choice of the \emph{compact} real form for the dynamics
(\autoref{post:app-albert:compact}) and the soldering of Lorentzian signature from
$N$ (\autoref{rmk:app-albert:signature}, postulate P3). The algebra fixes the
arena; the physics input is the potential class (P1) and the soldering (P3),
treated in their own units.
\end{remark}
